\section{Verbessert Aggregation wirklich die Detection}
Die zentrale Frage dieser Arbeit ist, ob die vorgestellte kollaborative Aggregation von CTI-Reports die Detektionsqualität der einzelnen Knoten empirisch messbar verbessert. Die folgenden Unterabschnitte identifizieren konkret Einflussgrößen (siehe Listing~\ref{lst:config}), deren systematische Untersuchung einer zukünftigen Arbeit unterliegt. 

\subsection{Datenqualität vs. Datenmenge}
Die implementierten Filtermechanismen \texttt{filterReportsWithNaN} und \texttt{filterReportsByTimeSkew} reduzieren die Anzahl eingehender CTI-Reports aktiv vor der Aggregation. Die zugrundeliegende These lautet: Ein vollständiger, zeitlich konsistenter Report trägt mehr zur Aggregationsqualität bei als eine größere Menge verrauschter oder unvollständiger Berichte. Es bleibt offen, ab welchem Filteranteil der verbleibende Informationsgehalt nicht mehr ausreicht, um eine zuverlässige Aggregation zu gewährleisten.

\subsection{Abtastintervall}
Wie in Kapitel~\ref{sec:treat_report_datenmanagement} (Threat-Report Datenmanagement) beschrieben, erfasst der CTI-Reporter aktuell drei Datenpunkte bei $t = 0$\,s, $t = -30$\,s und $t = -60$\,s. Die Wahl des 30-Sekunden-Intervalls stellt einen nicht empirisch validierten Kompromiss dar: Kürzere Intervalle erhöhen das Rauschen und die Auslastung im IPFS-Netzwerk, während längere Intervalle die Reaktionsfähigkeit auf laufende Angriffe verringern und veraltete Kontextinformation in die Aggregation einbringen. Der Einfluss variierender Intervalle auf die True-Positive-Rate des Gesamtsystems ist empirisch zu quantifizieren.

\subsection{Verschobene Zeitfenster}
Der Parameter \texttt{MAX\_TIME\_SKEW\_SECONDS} ist aktuell auf $90$\,s hardcodiert. In einer zukünftigen Arbeit muss untersucht werden, welcher Schwellenwert die Bedrohungseinschätzung durch den aggregierten CTI-Report nicht verzerrt und wie die Gesamtergebnisse dadurch optimiert werden können. 

\subsection{Echo-Chamber-Effekt}
Wie in Kapitel~\ref{sec:analyse_des_echo_chamber_effekts} (Analyse des Echo-Chamber-Effekts) formalisiert, bezeichnet der Echo-Chamber-Effekt das Szenario, in dem eine korrelierte Mehrheit kompromittierter Knoten den aggregierten CTI-Report systematisch verzerrt. Der definierte Grenzwert $W_{\text{max}} = 3$ dämpft den Einfluss einzelner Ausreißer, eliminiert jedoch nicht die strukturelle Verzerrung durch eine gleichartig manipulierte Mehrheit. Zur Adressierung dieser Forschungslücke kommen folgende Ansätze in Betracht: \todohl{die Integration von \textit{Byzantine Fault Tolerance}-Mechanismen in die Aggregationslogik, eine reputationsbasierte Knotengewichtung über die Zeit (\textit{Trust Score}) sowie die Kreuzvalidierung aggregierter Reports mit den lokalen Detektionsergebnissen des empfangenden Knotens als unabhängigem Anker.}{Quelle!!}

\subsection{Evaluation auf öffentlichem Datensatz}
Die Evaluation der vorliegenden Arbeit basiert ausschließlich auf einer physischen Testumgebung mit einer einzelnen ChargeIQ-Wallbox (siehe Kapitel~\ref{chapter:systemarchitektur} (Systemarchitektur und Detektion Framework)), was die externe Validität der Ergebnisse erheblich einschränkt. Ob das entwickelte Aggregations- und Gewichtungsverfahren auf Infrastrukturen außerhalb der eigenen Ladeumgebung generalisiert, bleibt ungeklärt. Als geeignete Evaluationsgrundlagen für zukünftige Arbeiten kommen etablierte IDS-Benchmarkdatensätze wie CICIDS2017/2018 \todohl{cite-sharafaldin2018cicids}{Datensatz?}, N-BaIoT \todohl{cite-meidan2018nbaiot}{Quelle bzw. geeigneten Datensatz finden} oder IoT-spezifische Datensätze in Betracht. Erst eine solche Evaluation ermöglicht eine wissenschaftlich belastbare Aussage über den tatsächlichen Mehrwert des kollaborativen Ansatzes gegenüber rein lokaler Anomalieerkennung.